\section{Sources of Bias}

\begin{frame}{First: ``Bias'' Means Three Different Things}

    \begin{columns}[T]
        \begin{column}{0.32\textwidth}
            \begin{tcolorbox}[colback=gray!8, colframe=gray!60, title=Statistical bias]
                \small
                $\mathbb{E}[\hat{\theta}] - \theta \neq 0$.

                \vspace{0.4em}
                The estimator is systematically off. Neutral term; half of the
                bias--variance trade-off.
            \end{tcolorbox}
        \end{column}
        \begin{column}{0.32\textwidth}
            \begin{tcolorbox}[colback=raiBlue!5, colframe=raiBlue!60, title=Inductive bias]
                \small
                The assumptions a model class makes.

                \vspace{0.4em}
                Convolutions assume locality; linear models assume additivity.
                \textbf{Necessary} --- without it there is no generalisation.
            \end{tcolorbox}
        \end{column}
        \begin{column}{0.32\textwidth}
            \begin{tcolorbox}[colback=raiRed!5, colframe=raiRed!70, title=Societal / harmful bias]
                \small
                Systematically different quality of service or outcome across groups of people.

                \vspace{0.4em}
                \textbf{This lecture's subject.}
            \end{tcolorbox}
        \end{column}
    \end{columns}

    \vspace{1.2em}

    \begin{block}{}
        \centering
        The three are unrelated. A model can be statistically unbiased and still harmful; a strong
        inductive bias can either mask or expose a societal one.
    \end{block}
\end{frame}
